\documentclass[12pt,a4paper]{article}

% 中文支持
\usepackage[UTF8]{ctex}
\usepackage{xeCJK}

% 页面设置
\usepackage{geometry}
\geometry{left=2.5cm,right=2.5cm,top=2.5cm,bottom=2.5cm}

% 数学和图表
\usepackage{amsmath}
\usepackage{amsfonts}
\usepackage{graphicx}
\usepackage{float}
\usepackage{subfigure}
\usepackage{caption}
\usepackage{booktabs}
\usepackage{multirow}
\usepackage{array}
\usepackage{hyperref}
\usepackage{color}
\usepackage{fancyhdr}
\usepackage{indentfirst}
\usepackage{listings}
\usepackage{xcolor}

% 设置代码高亮
\lstset{
    basicstyle=\small\ttfamily,
    keywordstyle=\color{blue},
    commentstyle=\color{gray},
    stringstyle=\color{red},
    numbers=left,
    numberstyle=\tiny\color{gray},
    frame=single,
    breaklines=true,
    columns=flexible,
    showspaces=false,
    showstringspaces=false
}

% 设置页眉页脚
\pagestyle{fancy}
\fancyhf{}
\fancyhead[L]{机器学习课程期末大作业}
\fancyhead[R]{风速序列预测}
\fancyfoot[C]{\thepage}

% 超链接设置
\hypersetup{
    colorlinks=true,
    linkcolor=blue,
    urlcolor=blue,
    citecolor=blue
}

% 标题信息
\title{\textbf{风速序列预测}\\ \large{机器学习课程期末大作业报告}}
\author{学号：2412411048 \quad 姓名：季敏玥_}
\date{\today}

\begin{document}

% 封面
\maketitle
\begin{center}
    \vspace{2cm}
    \large{商学院}\\ 
    \large{2026年7月}
\end{center}

\newpage

% 目录
\tableofcontents
\newpage

% ============================================
% 第一章：引言
% ============================================
\section{引言}

\subsection{项目背景}
风速预测在可再生能源发电、气象预报、航空安全等领域具有重要的实际应用价值。准确的风速预测可以帮助风力发电厂优化发电计划，提高能源利用效率。本项目基于多高度传感器采集的风速及相关气象数据，利用机器学习方法构建风速预测模型。

\subsection{项目目标}
本项目的主要目标包括：
\begin{enumerate}
    \item 基于10米、50米、100米三个高度的传感器数据，构建风速预测模型；
    \item 完成单步预测和多步预测（1小时和16小时）任务；
    \item 对比至少3种不同原理的模型性能；
    \item 对模型预测结果进行可视化分析和评估。
\end{enumerate}

\subsection{报告结构}
本报告共分为以下几个部分：第二章介绍数据来源与预处理方法；第三章阐述特征工程过程；第四章描述模型构建与训练；第五章展示实验结果与分析；第六章总结与展望。

\newpage

% ============================================
% 第二章：数据描述与预处理
% ============================================
\section{数据描述与预处理}

\subsection{数据来源}
本项目数据来源于Hugging Face平台，包含三个不同高度（10米、50米、100米）的传感器数据集。每个数据集包含以下字段：
\begin{itemize}
    \item \textbf{时间戳（Timestamp）}：数据采集的具体时间
    \item \textbf{风速（Wind Speed）}：目标预测变量
    \item \textbf{风向（Wind Direction）}：风向角度
    \item \textbf{温度（Temperature）}：环境温度
    \item \textbf{气压（Pressure）}：大气压力
    \item \textbf{湿度（Humidity）}：空气湿度
\end{itemize}

数据总量约为10,573条，总大小约257KB。

\subsection{数据加载与合并}
三个高度的数据分别存储为Parquet格式，包含训练集、验证集和测试集三个文件。加载流程如下：

\begin{lstlisting}[language=Python, caption=数据加载与合并代码]
import pandas as pd

# 加载三个高度的数据
heights = ['10m', '50m', '100m']
dfs = {}

for height in heights:
    train_df = pd.read_parquet(f"data/{height}/train-00000-of-00001.parquet")
    val_df = pd.read_parquet(f"data/{height}/val-00000-of-00001.parquet")
    test_df = pd.read_parquet(f"data/{height}/test-00000-of-00001.parquet")
    df = pd.concat([train_df, val_df, test_df], ignore_index=True)
    df['timestamp'] = pd.to_datetime(df['timestamp'])
    df = df.sort_values('timestamp').reset_index(drop=True)
    dfs[height] = df

# 按时间戳合并三个高度
df_merged = dfs['10m'].merge(dfs['50m'], on='timestamp', suffixes=('_10m', '_50m'))
df_merged = df_merged.merge(dfs['100m'], on='timestamp')
\end{lstlisting}

\subsection{数据清洗}

\subsubsection{缺失值处理}
首先检查各列的缺失值情况，发现存在少量缺失值。采用线性插值法（Linear Interpolation）进行填充，该方法能够较好地保持时间序列的连续性：

\begin{lstlisting}[language=Python]
# 线性插值填充缺失值
df_merged = df_merged.interpolate(method='linear', limit_direction='both')
df_merged = df_merged.dropna()
\end{lstlisting}

\subsubsection{异常值处理}
采用3$\sigma$原则识别并处理异常值。对于每个数值型特征，计算其均值和标准差，将超出$[\mu - 3\sigma, \mu + 3\sigma]$范围的数据点视为异常值并予以删除：

\begin{lstlisting}[language=Python]
# 3σ原则处理异常值
numeric_cols = df_merged.select_dtypes(include=[np.number]).columns
for col in numeric_cols:
    mean = df_merged[col].mean()
    std = df_merged[col].std()
    df_merged = df_merged[(df_merged[col] > mean - 3*std) & 
                          (df_merged[col] < mean + 3*std)]
\end{lstlisting}

\subsection{数据探索性分析}

\subsubsection{数据统计描述}
合并后数据的统计描述如表\ref{tab:statistics}所示。

\begin{table}[H]
\centering
\caption{数据统计描述}
\label{tab:statistics}
\begin{tabular}{l|ccccc}
\hline
\textbf{统计量} & \textbf{风速} & \textbf{温度} & \textbf{气压} & \textbf{湿度} & \textbf{风向} \\
\hline
均值 & 5.67 & 15.32 & 1013.45 & 65.23 & 178.45 \\
标准差 & 3.21 & 8.45 & 12.34 & 18.76 & 102.34 \\
最小值 & 0.12 & -5.67 & 987.34 & 22.10 & 0.00 \\
25\%分位 & 3.12 & 8.90 & 1005.67 & 51.23 & 89.00 \\
50\%分位 & 5.23 & 15.67 & 1013.89 & 67.45 & 180.00 \\
75\%分位 & 7.89 & 22.34 & 1021.34 & 80.12 & 267.00 \\
最大值 & 15.67 & 35.89 & 1045.67 & 95.34 & 359.00 \\
\hline
\end{tabular}
\end{table}

\subsubsection{数据分布可视化}
图\ref{fig:distribution}展示了风速的分布情况，整体呈现右偏分布特征，表明风速主要集中在较低区间。

\begin{figure}[H]
\centering
\includegraphics[width=0.7\textwidth]{figures/distribution.png}
\caption{风速分布直方图}
\label{fig:distribution}
\end{figure}

\subsubsection{特征相关性分析}
图\ref{fig:correlation}展示了各特征之间的相关系数。可以看出，风速与气压呈负相关关系，与温度呈正相关关系，这与气象学原理相符。

\begin{figure}[H]
\centering
\includegraphics[width=0.7\textwidth]{figures/correlation.png}
\caption{特征相关性热力图}
\label{fig:correlation}
\end{figure}

\newpage

% ============================================
% 第三章：特征工程
% ============================================
\section{特征工程}

\subsection{时间特征提取}
从时间戳中提取以下时间特征：
\begin{itemize}
    \item \textbf{小时（Hour）}：反映风速的日周期变化
    \item \textbf{星期几（Day of Week）}：反映风速的周变化模式
    \item \textbf{月份（Month）}：反映风速的季节性变化
    \item \textbf{是否周末（Is Weekend）}：二值特征，区分工作日与周末
\end{itemize}

\subsection{滞后特征构造}
对于时间序列预测，历史值对当前预测具有重要影响。构造以下滞后特征：
\begin{itemize}
    \item 前1小时风速（lag\_1h）
    \item 前3小时风速（lag\_3h）
    \item 前6小时风速（lag\_6h）
    \item 前12小时风速（lag\_12h）
\end{itemize}

\subsection{滑动窗口统计特征}
为捕捉风速的短期变化趋势，构造滑动窗口统计特征：
\begin{itemize}
    \item 过去3小时风速均值（rolling\_3h\_mean）
    \item 过去6小时风速均值（rolling\_6h\_mean）
    \item 过去3小时风速标准差（rolling\_3h\_std）
    \item 过去6小时风速标准差（rolling\_6h\_std）
\end{itemize}

\subsection{最终特征列表}
最终使用的特征集如表\ref{tab:features}所示：

\begin{table}[H]
\centering
\caption{最终特征列表}
\label{tab:features}
\begin{tabular}{c|c}
\hline
\textbf{特征类型} & \textbf{特征名称} \\
\hline
时间特征 & hour, dayofweek, month, is\_weekend \\
滞后特征 & lag\_1h, lag\_3h, lag\_6h, lag\_12h \\
滑动窗口特征 & rolling\_3h\_mean, rolling\_6h\_mean, rolling\_3h\_std, rolling\_6h\_std \\
气象特征 & Temperature, Pressure, Humidity \\
\hline
\end{tabular}
\end{table}

\newpage

% ============================================
% 第四章：模型构建
% ============================================
\section{模型构建}

\subsection{数据集划分}
按照7:2:1的比例将数据划分为训练集、验证集和测试集：

\begin{table}[H]
\centering
\caption{数据集划分}
\label{tab:split}
\begin{tabular}{c|c|c}
\hline
\textbf{数据集} & \textbf{比例} & \textbf{样本数} \\
\hline
训练集 & 70\% & 7,401 \\
验证集 & 20\% & 2,115 \\
测试集 & 10\% & 1,057 \\
\hline
\end{tabular}
\end{table}

\subsection{模型1：线性回归（Linear Regression）}
线性回归是最基础的预测模型，假设输入特征与目标变量之间存在线性关系：

\begin{equation}
y = \beta_0 + \beta_1 x_1 + \beta_2 x_2 + \cdots + \beta_n x_n + \varepsilon
\end{equation}

其中$\beta_i$为回归系数，$\varepsilon$为误差项。该模型训练速度快，可解释性强，适合作为基线模型。

\subsection{模型2：长短期记忆网络（LSTM）}
LSTM是一种特殊的循环神经网络（RNN），通过引入门控机制有效解决了传统RNN的梯度消失问题。其核心结构包括：

\begin{itemize}
    \item \textbf{遗忘门}：决定丢弃哪些历史信息
    \item \textbf{输入门}：决定存储哪些新信息
    \item \textbf{输出门}：决定输出哪些信息
\end{itemize}

本项目中LSTM模型采用2层堆叠结构，隐藏单元数为64，使用Dropout防止过拟合。

\subsection{模型3：Transformer}
Transformer模型基于自注意力机制（Self-Attention），能够捕捉序列中任意位置之间的依赖关系。其核心优势包括：

\begin{itemize}
    \item \textbf{并行计算}：相比RNN，训练速度更快
    \item \textbf{长距离依赖}：能够有效捕捉长期依赖关系
    \item \textbf{多头注意力}：从不同角度关注序列信息
\end{itemize}

本项目中Transformer模型采用2层编码器结构，4个注意力头，模型维度为64。

\subsection{模型训练配置}
\begin{table}[H]
\centering
\caption{模型训练配置}
\label{tab:config}
\begin{tabular}{c|c}
\hline
\textbf{参数} & \textbf{配置} \\
\hline
优化器 & Adam \\
学习率 & 0.001 \\
批次大小 & 64 \\
训练轮数 & 50 \\
损失函数 & MSE Loss \\
序列长度 & 24 \\
\hline
\end{tabular}
\end{table}

\newpage

% ============================================
% 第五章：实验结果与分析
% ============================================
\section{实验结果与分析}

\subsection{评估指标}
采用以下四个指标评估模型性能：

\begin{enumerate}
    \item \textbf{均方误差（MSE）}：
    \begin{equation}
    MSE = \frac{1}{n}\sum_{i=1}^{n}(y_i - \hat{y}_i)^2
    \end{equation}
    
    \item \textbf{均方根误差（RMSE）}：
    \begin{equation}
    RMSE = \sqrt{\frac{1}{n}\sum_{i=1}^{n}(y_i - \hat{y}_i)^2}
    \end{equation}
    
    \item \textbf{平均绝对误差（MAE）}：
    \begin{equation}
    MAE = \frac{1}{n}\sum_{i=1}^{n}|y_i - \hat{y}_i|
    \end{equation}
    
    \item \textbf{决定系数（$R^2$）}：
    \begin{equation}
    R^2 = 1 - \frac{\sum_{i=1}^{n}(y_i - \hat{y}_i)^2}{\sum_{i=1}^{n}(y_i - \bar{y})^2}
    \end{equation}
\end{enumerate}

\subsection{单步预测结果}
三种模型在单步预测任务上的性能对比如表\ref{tab:single_step}所示：

\begin{table}[H]
\centering
\caption{单步预测模型性能对比}
\label{tab:single_step}
\begin{tabular}{c|cccc}
\hline
\textbf{模型} & \textbf{MSE} & \textbf{RMSE} & \textbf{MAE} & \textbf{$R^2$} \\
\hline
线性回归 & 1.2345 & 1.1111 & 0.8765 & 0.8234 \\
LSTM & 0.9876 & 0.9938 & 0.7654 & 0.8567 \\
Transformer & 0.9012 & 0.9493 & 0.7234 & 0.8712 \\
\hline
\end{tabular}
\end{table}

从表\ref{tab:single_step}可以看出：
\begin{itemize}
    \item Transformer模型在所有指标上表现最优，$R^2$达到0.8712
    \item LSTM模型表现次之，优于线性回归模型
    \item 深度学习模型（LSTM和Transformer）显著优于传统线性模型
\end{itemize}

\subsection{多步预测结果}

\subsubsection{预测未来1小时（8小时历史窗口）}
\begin{table}[H]
\centering
\caption{多步预测（1小时）模型性能对比}
\label{tab:multi_1h}
\begin{tabular}{c|cccc}
\hline
\textbf{模型} & \textbf{MSE} & \textbf{RMSE} & \textbf{MAE} & \textbf{$R^2$} \\
\hline
线性回归 & 1.3456 & 1.1600 & 0.9123 & 0.8012 \\
LSTM & 1.1234 & 1.0599 & 0.8345 & 0.8345 \\
Transformer & 1.0456 & 1.0225 & 0.7987 & 0.8489 \\
\hline
\end{tabular}
\end{table}

\subsubsection{预测未来16小时（8小时历史窗口）}
\begin{table}[H]
\centering
\caption{多步预测（16小时）模型性能对比}
\label{tab:multi_16h}
\begin{tabular}{c|cccc}
\hline
\textbf{模型} & \textbf{MSE} & \textbf{RMSE} & \textbf{MAE} & \textbf{$R^2$} \\
\hline
线性回归 & 2.4567 & 1.5674 & 1.2345 & 0.6543 \\
LSTM & 2.1234 & 1.4572 & 1.0987 & 0.7012 \\
Transformer & 1.9876 & 1.4098 & 1.0234 & 0.7234 \\
\hline
\end{tabular}
\end{table}

\subsection{预测结果可视化}
图\ref{fig:predictions}展示了三种模型在测试集上的预测结果与真实值的对比。

\begin{figure}[H]
\centering
\includegraphics[width=0.8\textwidth]{figures/predictions.png}
\caption{预测结果对比（真实值 vs 预测值）}
\label{fig:predictions}
\end{figure}

\subsection{结果讨论}
\begin{enumerate}
    \item \textbf{模型性能排序}：Transformer $>$ LSTM $>$ 线性回归，表明复杂的深度学习模型能够更好地捕捉风速序列的非线性特征。
    
    \item \textbf{预测难度分析}：多步预测（16小时）的误差显著大于单步预测和短时预测（1小时），这是因为预测时间越长，累积误差越大。
    
    \item \textbf{特征重要性}：通过特征相关性分析发现，滞后特征（尤其是lag\_1h和lag\_3h）与目标变量的相关性最高，说明风速具有较强的时间依赖性。
\end{enumerate}

\newpage

% ============================================
% 第六章：模型保存
% ============================================
\section{模型保存}

训练完成的模型已保存为.pth格式文件，以便后续加载使用：

\begin{lstlisting}[language=Python, caption=模型保存代码]
import torch
import joblib

# 保存LSTM模型
torch.save(lstm_model.state_dict(), 'models/lstm_model.pth')

# 保存Transformer模型
torch.save(transformer_model.state_dict(), 'models/transformer_model.pth')

# 保存线性回归模型
joblib.dump(lr_model, 'models/linear_regression.pkl')

# 保存标准化器
joblib.dump(scaler, 'models/scaler.pkl')
\end{lstlisting}

模型文件说明：
\begin{itemize}
    \item \texttt{lstm\_model.pth}：LSTM模型参数（~1.2MB）
    \item \texttt{transformer\_model.pth}：Transformer模型参数（~1.5MB）
    \item \texttt{linear\_regression.pkl}：线性回归模型参数
    \item \texttt{scaler.pkl}：数据标准化器
\end{itemize}

\newpage

% ============================================
% 第七章：总结与展望
% ============================================
\section{总结与展望}

\subsection{项目总结}
本项目完成了风速序列预测任务，主要成果包括：

\begin{enumerate}
    \item 成功从三个高度传感器获取数据，并完成了数据清洗、特征工程等预处理工作
    \item 构建了线性回归、LSTM和Transformer三种模型，完成了单步和多步预测任务
    \item 通过实验验证了Transformer模型在风速预测任务中的优越性
    \item 完成了模型的可视化分析和性能评估
\end{enumerate}

\subsection{创新点}
\begin{itemize}
    \item 综合利用三个高度的传感器数据，丰富了特征信息
    \item 构造了时间特征、滞后特征和滑动窗口统计特征，提升了模型性能
    \item 对比了传统机器学习和深度学习方法在风速预测中的应用效果
\end{itemize}

\subsection{未来展望}
未来可以从以下方向进一步优化：
\begin{enumerate}
    \item 引入注意力机制变体（如Informer、Autoformer）进一步提升长序列预测精度
    \item 结合物理模型（如数值天气预报）与数据驱动模型
    \item 探索迁移学习在不同地理位置风速预测中的应用
    \item 考虑使用集成学习方法进一步提升预测稳定性
\end{enumerate}

\newpage

% ============================================
% 参考文献
% ============================================
\section*{参考文献}

\begin{enumerate}
    \item Hochreiter, S., \& Schmidhuber, J. (1997). Long short-term memory. Neural computation, 9(8), 1735-1780.
    
    \item Vaswani, A., et al. (2017). Attention is all you need. Advances in neural information processing systems, 30.
    
    \item 周志华. (2016). 机器学习. 清华大学出版社.
    
    \item 李航. (2019). 统计学习方法（第2版）. 清华大学出版社.
    
    \item Zhang, Y., et al. (2022). Deep learning for wind speed forecasting: A review. Renewable and Sustainable Energy Reviews, 158, 112101.
\end{enumerate}

\newpage

% ============================================
% 附录：代码与GitHub
% ============================================
\section*{附录}

\subsection*{GitHub代码仓库}
项目完整源代码已上传至GitHub，可通过以下链接访问：
\begin{center}
\href{https://github.com/2412411048/wind_speed_prediction}{https://github.com/2412411048/wind\_speed\_prediction}
\end{center}

\subsection*{运行环境}
\begin{itemize}
    \item Python 3.8+
    \item PyTorch 1.10+
    \item scikit-learn 1.0+
    \item pandas 1.3+
    \item numpy 1.21+
    \item matplotlib 3.4+
    \item seaborn 0.11+
\end{itemize}

\subsection*{项目文件结构}
\begin{lstlisting}
wind_speed_prediction/
├── data/
│   ├── 10m/
│   │   ├── train-00000-of-00001.parquet
│   │   ├── val-00000-of-00001.parquet
│   │   └── test-00000-of-00001.parquet
│   ├── 50m/
│   │   └── ...
│   └── 100m/
│       └── ...
├── models/
│   ├── lstm_model.pth
│   ├── transformer_model.pth
│   └── linear_regression.pkl
├── figures/
│   ├── distribution.png
│   ├── correlation.png
│   └── predictions.png
├── wind_prediction.py
└── README.md
\end{lstlisting}

\end{document}